% GENERATED FILE. Edit canonical lessons, then run npm run generate:books.
% canonical-source-hash: fnv1a64:8dcdcb7fa23b4ae4
% canonical-lessons: ES-C650-autocar, ES-C650-tranvia, ES-C650-bonobus, ES-C650-transbordo, ES-C650-validar

\chapter{Autocar y tranvía}
\label{ch:es-autocar-y-tranvia}
\hlchaptermodality{650}

\begin{chapteropening}
\textbf{By the end of this chapter:} \emph{I can take a coach or a tram, use a bus pass, change and validate my ticket.}

Complete the last lesson of chapter 650: I can take a coach or a tram, use a bus pass, change and validate my ticket.
\end{chapteropening}

\section[el autocar]{el autocar — the coach (long-distance bus)}

\label{lesson:ES-C650-autocar}



\begin{quote}
Before the new one: say the Spanish for to pedal, then the Spanish for the cyclist.
\end{quote}

\subsection*{You'll want to know: el autocar}
\textbf{el autocar} — \textquotedblleft{}the coach (long-distance bus)\textquotedblright{}.

\begin{quote}
El autocar sale a las ocho.
\end{quote}

\begin{quote}
\emph{The coach leaves at eight.}
\end{quote}

A Spain word; Latin America says \emph{el autobús}, \emph{el bus} or \emph{el micro}.

\begin{cousinweb}
\emph{autocar} was borrowed from French \emph{autocar} (EtymDB), made from \emph{auto-} and English \emph{car}.
\end{cousinweb}

\begin{grammarlens}[title={What you've built}]
One more word for telling, explaining and describing at B1.
\end{grammarlens}

\subsection*{Your turn}
\begin{itemize}
\raggedright
  \item \emph{Say it:} \emph{el autocar}
  \item \emph{Say it:} \emph{el autocar}, once more
  \item \emph{Say it:} a sentence of your own with \emph{el autocar}
  \item \emph{Recall:} say the Spanish for to pedal, then the Spanish for the cyclist, then say \emph{el autocar} again
\end{itemize}

\begin{tcolorbox}[breakable,colback=teal!4,colframe=teal!35!black,title={Before you move on}]
What does el autocar mean? (\textquotedblleft{}The coach (long-distance bus)\textquotedblright{}.) Say it once more.
\end{tcolorbox}
\section[el tranvía]{el tranvía — the tram, the streetcar}

\label{lesson:ES-C650-tranvia}



\begin{quote}
Before the new one: say the Spanish for the cyclist, then the Spanish for the coach.
\end{quote}

\subsection*{You'll want to know: el tranvía}
\textbf{el tranvía} — \textquotedblleft{}the tram, the streetcar\textquotedblright{}.

\begin{quote}
El tranvía pasa cada diez minutos.
\end{quote}

\begin{quote}
\emph{The tram comes every ten minutes.}
\end{quote}

It ends in -a but is masculine: \emph{el tranvía}.

\begin{cousinweb}
\emph{tranvía} was borrowed from English \emph{tramway}, with the second half remade as Spanish \emph{vía} \textquotedblleft{}track, way\textquotedblright{}, from Latin \emph{via}.
\end{cousinweb}

\begin{grammarlens}[title={What you've built}]
One more word for telling, explaining and describing at B1.
\end{grammarlens}

\subsection*{Your turn}
\begin{itemize}
\raggedright
  \item \emph{Say it:} \emph{el tranvía}
  \item \emph{Say it:} \emph{el tranvía}, once more
  \item \emph{Say it:} a sentence of your own with \emph{el tranvía}
  \item \emph{Recall:} say the Spanish for the cyclist, then the Spanish for the coach, then say \emph{el tranvía} again
\end{itemize}

\begin{tcolorbox}[breakable,colback=teal!4,colframe=teal!35!black,title={Before you move on}]
What does el tranvía mean? (\textquotedblleft{}The tram, the streetcar\textquotedblright{}.) Say it once more.
\end{tcolorbox}
\section[el bonobús]{el bonobús — the bus pass (for several trips)}

\label{lesson:ES-C650-bonobus}



\begin{quote}
Before the new one: say the Spanish for the coach, then the Spanish for the tram.
\end{quote}

\subsection*{You'll want to know: el bonobús}
\textbf{el bonobús} — \textquotedblleft{}the bus pass (for several trips)\textquotedblright{}.

\begin{quote}
Con el bonobús el viaje sale más barato.
\end{quote}

\begin{quote}
\emph{With a bus pass the trip works out cheaper.}
\end{quote}

A Spain word; it holds ten trips, or a month of travel.

\begin{cousinweb}
\emph{bonobús} is built from \emph{bono} \textquotedblleft{}voucher\textquotedblright{}, from Latin \emph{bonus} \textquotedblleft{}good\textquotedblright{}, and \emph{bus}, short for \emph{autobús}, from Latin \emph{omnibus} \textquotedblleft{}for all\textquotedblright{}.
\end{cousinweb}

\begin{grammarlens}[title={What you've built}]
One more word for telling, explaining and describing at B1.
\end{grammarlens}

\subsection*{Your turn}
\begin{itemize}
\raggedright
  \item \emph{Say it:} \emph{el bonobús}
  \item \emph{Say it:} \emph{el bonobús}, once more
  \item \emph{Say it:} a sentence of your own with \emph{el bonobús}
  \item \emph{Recall:} say the Spanish for the coach, then the Spanish for the tram, then say \emph{el bonobús} again
\end{itemize}

\begin{tcolorbox}[breakable,colback=teal!4,colframe=teal!35!black,title={Before you move on}]
What does el bonobús mean? (\textquotedblleft{}The bus pass (for several trips)\textquotedblright{}.) Say it once more.
\end{tcolorbox}
\section[el transbordo]{el transbordo — the change, the transfer (between trains or buses)}

\label{lesson:ES-C650-transbordo}



\begin{quote}
Before the new one: say the Spanish for the tram, then the Spanish for the bus pass.
\end{quote}

\subsection*{You'll want to know: el transbordo}
\textbf{el transbordo} — \textquotedblleft{}the change, the transfer (between trains or buses)\textquotedblright{}.

\begin{quote}
Hay que hacer transbordo en la siguiente estación.
\end{quote}

\begin{quote}
\emph{You have to change at the next station.}
\end{quote}

\emph{Hacer transbordo} — to change trains or lines.

\begin{cousinweb}
\emph{transbordo} is built on Latin \emph{trans} \textquotedblleft{}across\textquotedblright{} and \emph{bordo} \textquotedblleft{}side of a ship\textquotedblright{}, which comes from \emph{borde} \textquotedblleft{}edge\textquotedblright{} (EtymDB). It first meant moving goods from one ship to another.
\end{cousinweb}

\begin{grammarlens}[title={What you've built}]
One more word for telling, explaining and describing at B1.
\end{grammarlens}

\subsection*{Your turn}
\begin{itemize}
\raggedright
  \item \emph{Say it:} \emph{el transbordo}
  \item \emph{Say it:} \emph{el transbordo}, once more
  \item \emph{Say it:} a sentence of your own with \emph{el transbordo}
  \item \emph{Recall:} say the Spanish for the tram, then the Spanish for the bus pass, then say \emph{el transbordo} again
\end{itemize}

\begin{tcolorbox}[breakable,colback=teal!4,colframe=teal!35!black,title={Before you move on}]
What does el transbordo mean? (\textquotedblleft{}The change, the transfer (between trains or buses)\textquotedblright{}.) Say it once more.
\end{tcolorbox}
\section[validar]{validar — to validate (a ticket)}

\label{lesson:ES-C650-validar}



\begin{quote}
Before the new one: say the Spanish for the bus pass, then the Spanish for the change.
\end{quote}

\subsection*{You'll want to know: validar}
\textbf{validar} — \textquotedblleft{}to validate (a ticket)\textquotedblright{}.

\begin{quote}
Valide el billete al subir.
\end{quote}

\begin{quote}
\emph{Validate your ticket when you get on.}
\end{quote}

Buses and trams have a small machine for this.

\begin{cousinweb}
\emph{validar} is built on \emph{válido} \textquotedblleft{}valid\textquotedblright{}, a learned borrowing from Latin \emph{validus} \textquotedblleft{}strong, able\textquotedblright{} (Lewis \& Short), from \emph{valēre} \textquotedblleft{}to be strong\textquotedblright{}.
\end{cousinweb}

\begin{grammarlens}[title={What you've built}]
One more word for telling, explaining and describing at B1.
\end{grammarlens}

\subsection*{Your turn}
\begin{itemize}
\raggedright
  \item \emph{Say it:} \emph{validar}
  \item \emph{Say it:} \emph{validar}, once more
  \item \emph{Say it:} a sentence of your own with \emph{validar}
  \item \emph{Recall:} say the Spanish for the bus pass, then the Spanish for the change, then say \emph{validar} again
\end{itemize}

\begin{tcolorbox}[breakable,colback=teal!4,colframe=teal!35!black,title={Before you move on}]
What does validar mean? (\textquotedblleft{}To validate (a ticket)\textquotedblright{}.) Say it once more.
\end{tcolorbox}
